Figure~\ref{fig:atlas-box-convolution} interprets convolution as the length of overlap between two moving intervals.

\begin{figure}[!htbp]
\centering
\begin{tikzpicture}[>=Stealth]
\begin{axis}[width=5.8cm,height=4.8cm,axis lines=middle,ticklabel style={font=\small},label style={font=\small},legend style={font=\scriptsize,draw=black!20,fill=white},xlabel={$s$},xmin=-0.6,xmax=1.3,ymin=0,ymax=1.4,xtick={-0.25,0,0.75,1},xticklabels={$t-1$,$0$,$t$,$1$},ytick=\empty]
\addplot[figblue,very thick] coordinates {(0,1)(1,1)};
\addplot[figred,very thick] coordinates {(-0.25,0.65)(0.75,0.65)};
\addplot[figgreen,line width=4pt] coordinates {(0,0.3)(0.75,0.3)};
\draw[black!40,dashed] (axis cs:0,0) -- (axis cs:0,1.2);
\draw[black!40,dashed] (axis cs:0.75,0) -- (axis cs:0.75,1.2);
\end{axis}
\begin{axis}[width=5.8cm,height=4.8cm,axis lines=middle,ticklabel style={font=\small},label style={font=\small},legend style={font=\scriptsize,draw=black!20,fill=white},at={(6.2cm,0)},xlabel={$t$},ylabel={$(f*f)(t)$},xmin=-0.3,xmax=2.3,ymin=-0.1,ymax=1.4,xtick={0,1,2},ytick={1}]
\addplot[figblue,very thick] coordinates {(-0.3,0)(0,0)(1,1)(2,0)(2.3,0)};
\addplot[only marks,figgreen,mark=*,mark size=2.5pt] coordinates {(0.75,0.75)};
\draw[figgreen,dashed] (axis cs:0.75,0) -- (axis cs:0.75,0.75);
\end{axis}
\end{tikzpicture}
\caption{Left: at shift $t=3/4$, the intervals $[0,1]$ and $[t-1,t]$ overlap on a segment of length $3/4$. Right: as $t$ varies, that overlap length forms the triangular convolution of two unit-interval indicator functions. It is zero outside $[0,2]$ and reaches one at $t=1$.}
\label{fig:atlas-box-convolution}
\end{figure}